% ====================================================================
% jmp_beamer_preamble.tex --- shared preamble for the JMP seminar deck.
%
% REUSED VERBATIM from the companion theory talk
%   beamer/reference/Theory_talk/slides/Presentation.tex
%     "Jobs and Well-Being Measurement" (Haydar and Maniquet)
% so that the two decks read as one visual family.  The elements carried
% over unchanged are marked [THEORY-TALK] below; everything else is an
% addition this deck needs and the theory talk does not.
%
% Consistency is the point: same class options, same theme, same accent
% colour, same frame numbering, same appendix handling, same visible-link
% convention, same note-page template, same two-agent sketch palette.
% ====================================================================

% ---- [THEORY-TALK] class, theme and the deep-red accent -------------
\usetheme{metropolis}
\definecolor{deepred}{RGB}{150,25,35}
\setbeamercolor{frametitle}{bg=deepred}
\setbeamertemplate{navigation symbols}{}
% metropolis-defined accents (safe: these beamercolors exist in metropolis)
\setbeamercolor{progress bar}{fg=deepred}
\setbeamercolor{title separator}{fg=deepred}
% frame number as current/total, in deep red
\metroset{numbering=fraction}
\setbeamercolor{footline}{fg=deepred}
% backup appendix: keep the main frame total at the running-order count
\usepackage{appendixnumberbeamer}
% VISIBLE link colours for rehearsal (internal hyperlinks in deep red)
\hypersetup{colorlinks=true,linkcolor=deepred,citecolor=deepred,urlcolor=deepred}

% ---- [THEORY-TALK] encoding, language, maths, tables, tikz ----------
\usepackage[utf8]{inputenc}
\usepackage[english]{babel}
\usepackage{amsmath}
\usepackage{amssymb}
\usepackage{amsfonts}
\usepackage{booktabs}
\usepackage[official]{eurosym}   % \euro --- this deck quotes euro magnitudes
\usepackage{tikz}
\usetikzlibrary{decorations.pathreplacing,arrows.meta,positioning,fit,backgrounds}

% ---- [THEORY-TALK] the two-agent sketch palette --------------------
% red = agent i / household A ; blue = agent h / household B
\definecolor{agenti}{RGB}{200,30,40}
\definecolor{agenth}{RGB}{20,80,200}

% ---- [THEORY-TALK] notes on a second screen (rehearsal view) -------
% The rehearsal build turns this on; the projection build leaves it off.
\usepackage{pgfpages}
% Long notes overflow the note column by default; render them smaller,
% ragged-right, with tight paragraph spacing so the full text stays
% legible and inside the frame on the second screen.
%
% This deck's notes are longer than the theory talk's -- several run past
% one note page at \scriptsize and produced overfull \vbox warnings.  The
% template below keeps the theory talk's look and adds two things it
% needs: \tiny with a tighter baseline, and an \allowbreak-friendly
% vertical box so a long note breaks across note pages instead of
% overflowing.  Hyphenation is relaxed so no line runs into the margin.
\setbeamertemplate{note page}{%
  \vspace{0.3em}\par
  {\tiny\setlength{\parskip}{0.45em}\setlength{\parindent}{0pt}%
   \setlength{\emergencystretch}{3em}%
   \linespread{1.05}\selectfont
   \raggedright\insertnote\par}%
}
% Notes are prose, not code: let TeX hyphenate freely rather than push a
% long word past the measure.
\hyphenpenalty=50
\tolerance=2000

\DeclareMathOperator*{\argmax}{arg\,max}

% ====================================================================
% ADDITIONS for this deck (not in the theory talk)
% ====================================================================

% ---- channel colours ------------------------------------------------
% One colour per decomposition channel, used consistently on every slide
% that names a channel.  Preferences take the theory talk's agent-red so
% the two decks agree on which side of the argument is which.
\definecolor{chpref}{RGB}{200,30,40}     % preferences   (= agenti)
\definecolor{chacc}{RGB}{20,80,200}      % job access    (= agenth)
\definecolor{chearn}{RGB}{200,120,20}    % earning opportunities
\definecolor{chneeds}{RGB}{40,120,70}    % endowments and needs
\definecolor{chenv}{RGB}{60,60,70}       % the whole environment

% ---- the paper's notation ------------------------------------------
% Audience-facing channel names are WORDS on the slides; these commands
% exist so the words are spelled the same way every time they appear.
\newcommand{\pref}{\textcolor{chpref}{preferences}}
\newcommand{\Pref}{\textcolor{chpref}{Preferences}}
\newcommand{\acc}{\textcolor{chacc}{job access}}
\newcommand{\Acc}{\textcolor{chacc}{Job access}}
\newcommand{\earn}{\textcolor{chearn}{earning opportunities}}
\newcommand{\Earn}{\textcolor{chearn}{Earning opportunities}}
\newcommand{\needs}{\textcolor{chneeds}{household endowments and needs}}
\newcommand{\Needs}{\textcolor{chneeds}{Household endowments and needs}}
\newcommand{\env}{\textcolor{chenv}{the environment}}
\newcommand{\Env}{\textcolor{chenv}{The environment}}

% Symbols.  The theory-paper symbols W, z, R, A, y appear ONLY where the
% companion paper is cited; the estimation notation is separate.
\newcommand{\Wmeasure}{W}                      % theory-paper measure
\newcommand{\Jset}{\mathcal{J}}                % theory-paper job space
\newcommand{\Wone}{W^{1}}                      % theory-paper W^1
\newcommand{\yref}{\mathbf{y}^{w}}             % theory-paper reference pay
\newcommand{\gopp}{g}                          % estimated offer density
\newcommand{\uutil}{u}                         % preference index
\newcommand{\Wmm}{\mathrm{W}}                  % money-metric welfare level

% ---- headline / caveat call-outs -----------------------------------
\newcommand{\headline}[1]{{\large\textcolor{deepred}{#1}}}
\newcommand{\caveat}[1]{{\footnotesize\textcolor{deepred}{#1}}}
% a number with its numerical-integration band
\newcommand{\wband}[2]{#1\,\ensuremath{\pm}\,#2}

% ---- the 25-minute variant ------------------------------------------
% Every frame carries \shortdeck (in the 25-minute running order) or
% \longdeck (45-minute only).  \DeckShort, set by the driver file, drops
% the \longdeck frames; the default keeps everything.
%
% The frozen plan's 25-minute running order is 16 slides, reached from the
% 23 by FIVE cuts and TWO merges:
%   CUT   -> backup : 9, 12, 17, 18, 22          (\longdeck)
%   MERGE           : 6+7 -> one slide, 10+11 -> one slide
% So three tags, not two.  A merged pair is written as two full frames for
% the 45-minute deck; in the short deck the SECOND of the pair drops and
% the first shows its merged content.  23 - 5 - 2 = 16.
%
% Usage:  \shortdeck{\begin{frame}...\end{frame}}     kept in both
%         \longdeck{\begin{frame}...\end{frame}}      45-minute only (cut)
%         \mergedaway{\begin{frame}...\end{frame}}    45-minute only (merged)
%         \ifDeckShort ... \else ... \fi              inline merged content
\newif\ifDeckShort\DeckShortfalse
\newcommand{\shortdeck}[1]{#1}                       % always kept
\newcommand{\longdeck}[1]{\ifDeckShort\else#1\fi}    % cut to backup at 25 min
\newcommand{\mergedaway}[1]{\ifDeckShort\else#1\fi}  % folded into the previous
% content shown only in the merged (25-minute) form of a slide
\newcommand{\mergedonly}[1]{\ifDeckShort#1\fi}

% ---- figure helper ---------------------------------------------------
% Every figure is a PDF from the frozen sprint set, included by name from
% beamer/figures/.  \deckfig{name}{width fraction} keeps the call sites
% uniform and makes a missing file fail loudly at compile time.
\graphicspath{{figures/}}
% Heights leave room for the frametitle, the one caption line every figure
% slide carries, and the metropolis footline -- a taller box overflows the
% frame and shows up as an overfull \vbox.
\newcommand{\deckfig}[2]{%
  \begin{center}\includegraphics[width=#2\textwidth,height=0.70\textheight,%
    keepaspectratio]{#1}\end{center}}
\newcommand{\deckfigtwo}[4]{%
  \begin{columns}[c,onlytextwidth]
  \column{0.49\textwidth}\includegraphics[width=\textwidth,height=0.60\textheight,%
    keepaspectratio]{#1}\\{\centering\scriptsize #2\par}
  \column{0.49\textwidth}\includegraphics[width=\textwidth,height=0.60\textheight,%
    keepaspectratio]{#3}\\{\centering\scriptsize #4\par}
  \end{columns}}

% ---- tables ---------------------------------------------------------
% Several slides carry a comparison table that is naturally wider than the
% text block.  \decktable shrinks one to the measure rather than letting it
% run into the margin (an overfull \hbox), and never enlarges a narrow one.
\newcommand{\decktable}[1]{%
  \begin{center}\resizebox{\textwidth}{!}{#1}\end{center}}

\endinput
